\begin{examplebox}

	\textbf{\textcolor{CExample}{例 1（基础）}}

	\textbf{\textcolor{CExample}{题目：}}
	已知等差数列$\{a_n\}$中，$a_1=1$，$d=2$，取$m=3$。求$T_1$、$T_2$、$T_3$，并验证它们构成等差数列。

	\textbf{\textcolor{CExample}{解题步骤：}}
	\begin{description}[style=nextline, leftmargin=2.8em, labelsep=0.5em, itemsep=0.45em]
		\item[\textbf{第一步（求$S_3$）：}] 计算前$3$项和：
			\[
				\begin{aligned}
					S_3 &= 1+3+5 \\
					&= 9
			\end{aligned}
			\]
			或
			\[
				\begin{aligned}
					S_3 &= 3\times1 + \frac{3\times2}{2}\times2 \\
					&= 9.
			\end{aligned}
			\]
			\par\textbf{理由：} 等差数列求和公式或直接相加。

		\item[\textbf{第二步（计算$T_1,T_2,T_3$）：}] 根据定义：
			\[
				\begin{aligned}
					T_1 &= S_3 \\
					&= 9, \\[2pt]
					T_2 &= S_6 - S_3 \\
					&= (1+3+5+7+9+11)-9 \\
					&= 36-9 \\
					&= 27, \\[2pt]
					T_3 &= S_9 - S_6 \\
					&= (1+3+\cdots+17)-36 \\
					&= 81-36 \\
					&= 45.
			\end{aligned}
			\]
			\par\textbf{理由:} $T_k = S_{3k} - S_{3(k-1)}$,且$S_0=0$.

		\item[\textbf{第三步(验证等差):}] 计算相邻差:
			\[
				\begin{aligned}
					T_2 - T_1 &= 27 - 9 \\
					&= 18, \\
					T_3 - T_2 &= 45 - 27 \\
					&= 18.
			\end{aligned}
			\]
			\par\textbf{理由:} 两差相等,故$\{T_1,T_2,T_3\}$是公差为$18$的等差数列,且$m^2 d = 9\times2 = 18$,与结论一致.
	\end{description}

	\textbf{\textcolor{CExample}{关键结论:}}
	$T_1,T_2,T_3$成等差数列,公差$\boldsymbol{18}$.

	\medskip
	\textbf{\textcolor{CExample}{例 2(稍变形)}}

	\textbf{\textcolor{CExample}{题目:}}
	已知等差数列$\{a_n\}$的前$n$项和$S_n = 2n^2 + n$,取$m=2$,求第$5$个片段和$T_5$.

	\textbf{\textcolor{CExample}{解题步骤:}}
	\begin{description}[style=nextline, leftmargin=2.8em, labelsep=0.5em, itemsep=0.45em]
		\item[\textbf{第一步(求原公差$d$与首项$a_1$):}] 由$S_n$公式对比标准形式$S_n = n a_1 + \frac{n(n-1)}{2}d$,得到:
			\[
				\frac{d}{2} = 2,\quad a_1 - \frac{d}{2} = 1 \;\Rightarrow\; d=4,\; a_1 = 3.
			\]
			\par\textbf{理由:} 二次项系数匹配得$d$,一次项系数匹配得$a_1$.

		\item[\textbf{第二步(计算$S_m$与$m^2 d$):}] 已知$m=2$,则
			\[
				\begin{aligned}
					S_2 &= 2\cdot 2^2 + 2 \\
					&= 10, \\[2pt]
					m^2 d &= 2^2 \times 4 \\
					&= 16.
			\end{aligned}
			\]
			\par\textbf{理由:} 直接代入$S_n$公式或利用$a_1,d$计算$S_2$,并计算$m^2 d$.

		\item[\textbf{第三步(利用通项公式求$T_5$):}] 由结论通项公式$T_k = S_m + (k-1)m^2 d$,取$k=5$:
			\[
				\begin{aligned}
					T_5 &= 10 + (5-1)\times 16 \\
					&= 10 + 64 \\
					&= 74.
			\end{aligned}
			\]
			\par\textbf{理由:} 将已知数值代入通项公式即得结果.
	\end{description}

	\textbf{\textcolor{CExample}{关键结论:}}
	$\boldsymbol{T_5 = 74}$.

	\medskip
	\textbf{\textcolor{CExample}{例 3(综合应用)}}

	\textbf{\textcolor{CExample}{题目:}}
	在等差数列$\{a_n\}$中,已知$S_3=9$,$S_9=45$,求$S_6$.

	\textbf{\textcolor{CExample}{解题步骤:}}
	\begin{description}[style=nextline, leftmargin=2.8em, labelsep=0.5em, itemsep=0.45em]
		\item[\textbf{第一步(识别分段长度):}] 观察下标$3,6,9$,它们是$3$的倍数,因此$m=3$,$T_1=S_3$,$T_2=S_6-S_3$,$T_3=S_9-S_6$成等差数列.
			\[
				T_1=S_3,\quad T_2=S_6-S_3,\quad T_3=S_9-S_6.
			\]
			\par\textbf{理由:} 由结论知,对等差数列按固定$m$分段后,各片段和构成等差数列,这里$m=3$.

		\item[\textbf{第二步(列等差中项方程):}] 利用等差中项$2T_2 = T_1 + T_3$,得到方程:
			\[
				2(S_6 - S_3) = S_3 + (S_9 - S_6).
			\]
			\par\textbf{理由:} 等差数列$\{T_k\}$中,$T_1,T_2,T_3$满足$2T_2 = T_1 + T_3$.

		\item[\textbf{第三步(解方程得结果):}] 将已知$S_3=9$,$S_9=45$代入并求解:
			\[
				\begin{aligned}
					2(S_6 - 9) &= 9 + (45 - S_6) \\
					\Rightarrow\; 2S_6 - 18 &= 54 - S_6 \\
					\Rightarrow\; 3S_6 &= 72 \\
					\Rightarrow\; S_6 &= 24.
			\end{aligned}
			\]
			\par\textbf{理由:} 移项合并,最后除以系数得到$S_6$.
	\end{description}

	\textbf{\textcolor{CExample}{关键结论:}}
	$\boldsymbol{S_6 = 24}$.

\end{examplebox}
